\chapter{Related Work}

% ============================================================
% CONTENT MAP -- this chapter IS the literature review the CDS
% structure requires (PLAN.md rule 9); it's just named "Related Work"
% in the handbook's own list. The sources here are already real and
% peer-reviewed (docs/CITATIONS.md, 33 entries now in
% attachments/bibliography.bib). The work remaining is synthesis:
% group and compare, don't just list (PLAN.md rule 1).
%
%   \section{Indirect Prompt Injection}
%     \cite{greshake2023not} (foundational -- establishes the attack class),
%     \cite{liu2024formalizing} (formalization/benchmarking methodology --
%     compare your ASR methodology against this explicitly),
%     \cite{zhan2024injecagent} (tool-integrated agents -- adjacent/context
%     work, not directly used, say so)
%
%   \section{Knowledge Poisoning of RAG Systems}
%     \cite{zou2025poisonedrag} (the attack this thesis implements),
%     \cite{moradi2026defending} (2026 defense benchmarked against
%     PoisonedRAG -- evidence the attack remains the field's active
%     reference point)
%
%   \section{Multi-Turn Jailbreaking}
%     \cite{russinovich2025great} (Crescendo, implemented directly),
%     \cite{chao2024jailbreakbench}, \cite{mazeika2024harmbench}
%     (sources of the stratified behavior pool -- 20\%/80\% split),
%     \cite{nakka2026bitbypass} (contemporaneous 2026 work)
%
%   \section{Defenses}
%     \cite{hines2024defending} (spotlighting, implemented directly),
%     \cite{he2023debertav3} (architecture behind the instruction-
%     detection classifier), \cite{liu2025datasentinel} (closest
%     peer-reviewed comparator -- explain why a DIFFERENT classifier,
%     protectai/deberta-v3-base-prompt-injection-v2, was used in
%     practice), \cite{metaai2025llamaguard4} (output filter -- cite as
%     a released model, not a study), \cite{cui2025orbench},
%     \cite{rottger2024xstest} (over-refusal / false-positive
%     benchmarking used in Phase 3)
%
%   \section{Gap This Thesis Addresses}
%     - no existing work evaluates all three attack families against
%       the SAME four models under a SHARED, backend-controlled harness
%       (the backend-confound discovery itself, Approach/Evaluation
%       chapters, is a methodological point most of this literature
%       does not control for -- worth saying explicitly here)
%     - no existing work connects empirical RAG robustness results to
%       a concrete EU AI Act Article 15 compliance-testing gap
%       (PHASE4_EU_AI_ACT_MAPPING.md Sec.4's four contribution points --
%       preview here, full mapping goes in Evaluation/Conclusion)
%
% Suggested word target: ~3,200 words
% ============================================================
